% !TeX root = ../tcc.tex
\begin{resumo}[Abstract]
\begin{otherlanguage*}{english}
Leprosy is a neglected tropical disease whose diagnosis relies mainly on
clinical assessment, and the similarity of its lesions to other
dermatological conditions hinders early detection. Computer vision models
have been proposed to support the identification of these lesions, but the
way their performance is evaluated remains little examined. This work
analyzes the performance evaluation of computer vision models applied to the
identification of patterns associated with leprosy in skin images and, based
on the evidence gathered, plans a data science study. A structured
literature review, adapted from the systematic review protocol and with a
search assisted by artificial intelligence tools, included 26 studies, 17 of
which present their own data. Convolutional networks trained by transfer
learning and evaluation on a single split of the study's own dataset
prevail: of the 16 studies that develop and evaluate models, 14 use such a
split, only 2 report separating images by patient, and only 1 tests the
model on another dataset. Accuracy and sensitivity are the most reported
metrics, no study reports performance by skin tone, and differences in
datasets, classes, and procedures prevent a numerical comparison of the
results. Based on these gaps, a study was planned following the ACM SIGSOFT
Data Science standard, which will compare ResNet-50 and MobileNetV2 in the
classification of clinical photographs from the public 32 CSD v2
collection, with partitions by groups of related photographs, model
selection by cross-validation, balanced accuracy as the primary measure,
bootstrap confidence intervals, and metadata baselines. The study will be
carried out in the second stage of this monograph.

\vspace{\onelineskip}

\noindent
\textbf{Key-words}: leprosy. computer vision. performance evaluation. deep
learning. structured literature review. data science.
\end{otherlanguage*}
\end{resumo}
